\twocolumn[{
\begin{@twocolumnfalse}
\begin{center}
    {\LARGE\bfseries A Configurable Virtual Reality Sandbox for Comparing Prism-Effect Implementations}\\[0.5em]
    {\large Markus Birch Flensborg, Mathias Bisgaard, Rikke Bragh Jensen, Tze Huo Gucci Ho}\\[0.25em]
    {\normalsize Aalborg University, Medialogy}
\end{center}

\begin{abstract}
Prism adaptation is a rehabilitation method for unilateral spatial neglect in which repeated pointing under a visual displacement can produce error reduction during exposure and an after-effect after the displacement is removed. Virtual reality makes it possible to simulate prism-like perturbations in software, but current VR implementations differ in what is shifted, how feedback is shown, and which tasks are used for assessment. This makes comparison across studies difficult. We present a Unity-based VR sandbox that places translation, rotation, and skew perturbations inside one shared baseline--exposure--post workflow. The prototype supports controller-based and embodied hand-tracking input, logs event and movement data, and includes an RShiny dashboard for condition, participant, quality, spatial, and trajectory analysis. A healthy-participant evaluation was conducted with 10 participants and 80 completed modules. The system produced analyzable data across tasks and perturbations, but the measured after-effects were noisy and the no-effect condition was often non-zero. The clearest condition-level pattern was found for controller-based line bisection, where perturbation conditions showed positive median shifts above participant-specific no-effect baselines. Simulator Sickness Questionnaire responses showed a small post-session increase in discomfort, but session quality was more strongly limited by behaviour, posture, tracking, fatigue, and carry-over. The study indicates that a unified framework is useful not only for comparing perturbation types, but also for exposing procedural variables such as feedback visibility, confirmation method, tracking stability, posture, and carry-over between modules.
\end{abstract}
\vspace{1em}
\end{@twocolumnfalse}
}]
